\documentclass{article}
\usepackage{log_2022}           

\usepackage{booktabs}            % professional-quality tables
\usepackage{multirow}            % tabular cells spanning multiple rows
\usepackage{amsfonts}            % blackboard math symbols
\usepackage{graphicx}            % figures
\usepackage{duckuments}          % sample images

% If you want to use natbib:
\usepackage[numbers,compress,sort]{natbib}
%                                % for numerical citations
% \usepackage[sort,round]{natbib}
%                                % for textual citations

% If you want to use bibLaTeX, uncomment statements below:
% \usepackage[
%      backend=biber,
%      style=numeric-comp,
%      backref=true,
%      natbib=true]{biblatex}
% \addbibresource{reference.bib}

\title[MoML example style]{MoML example style}

\author[Y. Zhu et al.]{%
Yanqiao Zhu\thanks{Equal contribution.}\\
\institute{University of California, Los Angeles}\\
\email{yzhu@cs.ucla.edu}\And
Yuanqi Du\footnotemark[1]\\
\institute{Cornell University}\\
\email{yd392@cornell.edu}
}


\begin{document}

\maketitle

\begin{abstract}
Abstracts should be a single paragraph, ideally between 4--6 sentences long.
\end{abstract}

\section{Submission Guidelines}

We welcome papers related to machine learning for molecules, from small drug-like molecules to proteins and organic crystals. Papers can be submitted on the MoML website.




With the provided style file, the author information can be set in various styles.
If all authors are from the same institution, authors can use:
\begin{quote}
	\begin{verbatim}
		\author[F. Last et al.]
		 {First Last 1, First Last 2, ... \and First Last n \\
		  \institute{Institute \\ Address line}\\
		  \email{first.last@example.com}}
	\end{verbatim}
\end{quote}
For authors from different institutions, please use the \verb+\And+ command:
\begin{quote}
	\begin{verbatim}
		\author[F. Last et al.]
		 {First Last 1 \\
		  \institute{Institute 1\\ Address line} \\
		  \email{first.last.1@example.com}\And
		  First Last 2 \\
		  \institute{Institute 2\\ Address line} \\
		  \email{first.last.2@example.com}}
	\end{verbatim}
\end{quote}
If author names do not fit in one line, use the \verb+\AND+ command to start a separate row of authors:
\begin{quote}
	\begin{verbatim}
		\author[F. Last et al.]
		 {First Last 1, ..., \and First Last n \\
		  \institute{Institute 1\\ Address line} \\
		  \email{first.last.1@example.com}\AND
		  First Last n+1 \\
		  \institute{Institute 2\\ Address line} \\
		  \email{first.last.n.1@example.com}}
	\end{verbatim}
\end{quote}

\subsection{Figures}

\begin{figure}
	\centering
	\includegraphics[width=0.5\linewidth]{example-image-duck}
	\caption{A sample figure.}
	\label{fig:sample}
\end{figure}

All included artwork must be neat, legible, and separated from the text.
The figure number and caption always appear below the figure.
The figure label should be in boldface and numbered consecutively.
The caption should be set in 9-point type, in sentence case, and centered unless it runs more than one lines, in which case it should be flush left.
See Figure~\ref{fig:sample} for an example.

\subsection{Tables}

\begin{table}
	\centering
	\caption{A sample table.}
	\label{tab:sample}
	\begin{tabular}{ccccc}
		\toprule
		\multirow{2.5}{*}{Method} & \multicolumn{2}{c}{Data 1} & \multicolumn{2}{c}{Data 2}  \\
		\cmidrule(lr){2-3} \cmidrule(lr){4-5}
		& \(\mathbf{X}\) & \(\mathbf{Y}\) & \(\mathbf{X}\) & \(\mathbf{Y}\) \\
		\midrule
		A & 0.8817  & 0.9572 & 0.1893 & 0.1725 \\
		B & 0.7126  & 0.2615 & 0.9173 & 0.1286 \\
		C & 0.2716  & 0.1826 & 0.2836 & 0.1836 \\
		\bottomrule
	\end{tabular}
\end{table}

Like figures, tables should be legible and numbered consecutively.
However, the table number and caption should always appear above the table.
See Table~\ref{tab:sample} for an example.

Note that publication-quality tables \emph{do not contain vertical rules and double rules}.
We strongly suggest the use of the \texttt{booktabs} package which provides the commands \verb+\toprule+, \verb+\midrule+, and \verb+\bottomrule+ to enhance the quality of tables.

\subsection{Paragraphs}

Paragraphs are separated by 1/2 line space (5.5 points).
Do not indent the first line of a given paragraph.

\subsection{Equations}

The provided style file loads the \verb+amsmath+ package automatically.
Unnumbered single-lined equations should be displayed using \verb+\[+ and \verb+\]+. For example:
\[
	\mathbf{X}' = \sigma(\widetilde{\mathbf{D}}^{-\frac{1}{2}}\widetilde{\mathbf{A}}\widetilde{\mathbf{D}}^{-\frac{1}{2}} \mathbf{XW}).
\]
Numbered single-line equations should be displayed using the \verb+equation+ environment. For example:
\begin{equation}
	\mathbf{X}' = \sigma(\widetilde{\mathbf{D}}^{-1}\widetilde{\mathbf{A}}\mathbf{XW}).
\end{equation}

\subsection{Bibliographies}

Use an unnumbered first-level heading for the references.
For a citation, use \verb+\cite+, e.g.,~\cite{Kipf:2017tc}.
For a textual citation, use \verb+\citet+, e.g.,~\citet{Velickovic:2018we}.
\emph{Any choice} of citation style is allowed as long as it is used consistently throughout the whole paper.
Additionally, both \verb+natbib+ and \verb+bibLaTeX+ packages are supported.
It is also possible to reduce the font size to \verb+\small+ (9-point font) when listing the references.


\section*{Acknowledgements}

The \LaTeX{} template is heavily borrowed from LoG 2022.

The acknowledgements do not count towards the page limit.

% For natbib users:
\bibliographystyle{unsrtnat}
\bibliography{reference}
% For bibLaTeX users:
% \printbibliography

\appendix
\section{Appendix}
Any possible appendices should be placed after bibliographies.
If your paper has appendices, please submit the appendices together with the main body of the paper.
There will be no separate supplementary material submission.
The main text should be self-contained; reviewers are not obliged to look at the appendices when writing their review comments.

\end{document}
